% Generated draft: requires mathematical and source review.
\documentclass[11pt]{article}
\usepackage[margin=25mm]{geometry}
\usepackage{fontspec}
\setmainfont{Latin Modern Roman}
\usepackage{amsmath,amssymb,mathtools,unicode-math}
\setmathfont{Latin Modern Math}
\usepackage{enumitem,needspace,xurl,hyperref}
\setlength{\parindent}{0pt}
\setlength{\parskip}{5pt}
\setlength{\emergencystretch}{3em}
\newcommand{\N}{\mathbb{N}}
\newcommand{\Z}{\mathbb{Z}}
\newcommand{\Q}{\mathbb{Q}}
\newcommand{\R}{\mathbb{R}}
\newcommand{\C}{\mathbb{C}}
\newcommand{\F}{\mathbb{F}}
\newcommand{\A}{\mathbb{A}}
\newcommand{\PP}{\mathbb{P}}
\newcommand{\E}{\mathbb{E}}
\newcommand{\eps}{\varepsilon}
\newcommand{\dd}{\,\mathrm{d}}
\begin{document}
\section*{UnsolvedMath ID: 20000217}
\textbf{Source problem code:} catalogue entry
\subsection{Definitions and mathematical statement}
This extracted record is contextual theorem text, not a standalone open question. In the setting of a flag manifold $G/P$, let $W/W_P$ index Schubert varieties $X_s$. The source calls $X\subset G/P$ multiplicity-free when its cohomology class has Schubert coefficients zero or one:
\[
 [X]=\sum_{s\in S}[X_s],\qquad S\subset W/W_P,
\]
with the summands of the appropriate dimension. Equivalently in the intended intersection-theoretic setting, each zero-dimensional intersection with a general translate of a complementary Schubert variety is empty or one point. The malformed extracted summation index is interpreted this way, as its accompanying prose specifies.

Theorem 2.1 in the supplied source, attributed to Brion, states that a multiplicity-free $X$ degenerates to the reduced union $\bigcup_{s\in S}X_s$, and that this union is Cohen--Macaulay. The source further concludes that a reduced irreducible multiplicity-free $X$ is Cohen--Macaulay. Here a degeneration is understood as a flat family with the stated special fiber and general fiber isomorphic to $X$ in the appropriate embedded sense. The field, exact group hypotheses, and conventions in the word variety are not supplied; an arbitrary possibly nonreduced or reducible scheme is not silently included. The supplied verification says the separate classification question starts in the next canonical record.

The accompanying research report adds a structural claim for an integral $d$-dimensional multiplicity-free variety $X\subset(\PP^1)^n$. Its coordinate-projection matroid has bases $B$ of size $d$ for which $X\to(\PP^1)^B$ is dominant. The report claims every such projection is birational, so in affine coordinate functions
\[
 k(X)=k(t_b:b\in B).
\]
For adjacent bases $B=C\cup\{i\}$ and $B'=C\cup\{j\}$, it claims
\[
 t_j=\frac{a t_i+b}{c t_i+e},\qquad
 a,b,c,e\in k(t_c:c\in C),\qquad ae-bc\ne0.
\]
Thus basis exchange is a relative $\operatorname{PGL}_2$ transformation over the common coordinate field, not necessarily over the constant field. The induced birational basis-chart transitions satisfy the cycle cocycle: their composition around any closed path in the basis graph is the identity, with the requisite substitutions between coordinate fields.

The report also gives a rational-function realization reduction: seek coordinate functions in $k(z_1,\ldots,z_d)$ having the prescribed projection matroid, such that every matroid basis generates this whole field. Conversely their image closure in $(\PP^1)^n$ gives the proposed multiplicity-free realization under the intended multidegree conventions. Specifying the prescribed matroid is essential; basis-field equalities alone do not exclude additional bases. This is an asserted structural reduction, not a combinatorial classification, and its exact field and multidegree hypotheses require review.
\subsection{Short English statement}
Record Brion's stated multiplicity-free degeneration theorem and its Cohen--Macaulay consequence, rather than invent an open question from this contextual excerpt. The supplied additional report describes birational basis projections and relative fractional-linear exchange maps for multiplicity-free subvarieties of $(\PP^1)^n$, with a rational-function realization reduction but no combinatorial classification.
\subsection{Sources}
\begin{itemize}
\item Local registry: \url{data/unsolvedmath/problems.json}, permanent ID $20000217$, source code \texttt{catalogue entry}; extracted theorem text and full supplied research summary.
\item AIM workshop ``Degenerations in algebraic geometry'', problem-session notes, item 2 suggested by Allen Knutson, Theorem 2.1 attributed to Brion: \url{https://aimath.org/pastworkshops/degenalggeomproblems.pdf}. Supplied verification describes a six-page PDF dated September 7, 2016; it was not opened here. Canonical location \texttt{aim-algebraic-geometry-notes.json}, \texttt{notes[216]}.
\item No full bibliographic reference for Brion is supplied; original remarks and literature field empty. The supplied verification distinguishes this contextual record from the separately numbered classification question following it.
\item Supplied machine-generated research attempt 1: context only; problem status at run ``not a problem'', registry status open and classification \texttt{OPEN-TRIAGE}; rigor metadata ``proved'', novelty confidence low. Its additional structural claims are non-peer-reviewed and require independent verification before citation.
\end{itemize}
\end{document}
